\scanpage{267}
\mspage{781}

\unsure{mant|me} par la commutation des diagrammes
\[
(107)\quad
\diagram{\begin{tikzcd}[ampersand replacement=\&, column sep=large,row sep=small]
H(A) \arrow[dd,"(100)"',"\sim"] \arrow[rd,"\substack{i_A\\\sim}"] \& \\
\& H(\unsure{\mathbb{S}}_x^{\sim},\unsure{\mathbb{D}}_x^{*\sim}) \\
H(A') \arrow[ru,"\sim","i_{A'}"'] \&
\end{tikzcd}}
\;.
\]

Pour vérifier le fait que le foncteur $\underline{\mathrm{Gr}}(H)$ de (102)
est \unsure{ess\textsuperscript{t}.} \unsure{const\textsuperscript{t}},
je considère les deux foncteurs (indépendants du choix de
$\unsure{\mathbb{S}}_x^{\sim},\unsure{\mathbb{D}}_x^{*\sim}$)
\[
(108)\quad
\left\{
\diagram{\begin{tikzcd}[ampersand replacement=\&, column sep=small,row sep=large]
A \arrow[r,mapsto] \& \unsure{\mathbb{S}}_{xA}^{\sim},
\& A \arrow[r,mapsto] \& \unsure{\mathbb{D}}_{x,A}^{*\sim} \\
\underline{\mathcal{A}}_x \arrow[r] \& \underline{\unsure{\mathrm{Revuni}}}(\unsure{\mathbb{S}}_x)
\& \underline{\mathcal{A}}_x \arrow[r] \& \underline{\unsure{\mathrm{Revuni}}}(\unsure{\mathbb{D}}_x^*)
\end{tikzcd}}
\right.
\;,
\]
qui se factorisent en des \uline{équivalences}
\[
(109)\quad
\diagram{\begin{tikzcd}[ampersand replacement=\&, column sep=small]
\underline{\mathrm{Gr}}(\underline{\mathcal{A}}_x)
\arrow[r,"\substack{F\\\approx}"] \&
\underline{\unsure{\mathrm{Revuni}}}(\unsure{\mathbb{S}}_x),
\& \underline{\mathrm{Gr}}(\underline{\mathcal{A}}_x)
\arrow[r,"\substack{G\\\approx}"] \&
\underline{\unsure{\mathrm{Revuni}}}(\unsure{\mathbb{D}}_x^*)
\end{tikzcd}}
\]
Ceci dit, \inserted{utilisant
$\unsure{\mathbb{S}}_x^{\sim},\unsure{\mathbb{D}}_x^{*\sim}$,}
on a pour $A\in\mathrm{Ob}\,\underline{\mathcal{A}}_x=\mathcal{A}_x$ aussi,
des isom. fonctoriels en $A$
\[
(110)\quad\left\{
\begin{array}{l}
\unsure{\mathbb{S}^{\sim}|H}(A)
\simeq\mathrm{Isom}_{\underline{\unsure{\mathrm{Revuni}}}(\unsure{\mathbb{S}}_x)}
\deleted{\unsure{\mathbb{S}}}
\left(\underbrace{F(A)}_{\substack{\Vert\\\unsure{\mathbb{S}}_{x,A}^{\sim}}},\unsure{\mathbb{S}}^{\sim}\right)
\\[1.8em]
\unsure{\mathbb{D}}_x^{*\sim}(\unsure{\mathrm{Ep}|E_p}A)
\simeq\mathrm{Isom}_{\underline{\unsure{\mathrm{Revuni}}}(\unsure{\mathbb{D}}_x^*)}
\left(\underbrace{G(A)}_{\substack{\Vert\\\unsure{\mathbb{D}}_{x,A}^{*\sim}}},\unsure{\mathbb{D}}_x^{*\sim}\right),
\end{array}\right.
\]
et par la propriété univ. de $\underline{\mathrm{Gr}}(\underline{\mathcal{A}}_x)$,
ces isom. fonctoriels s'étendent sur la catégorie
$\underline{\mathrm{Gr}}(\underline{\mathcal{A}}_x)$.
Donc la situation est celle où on a un groupe comm. $\unsure{T|\overline T}$ $(=T_x)$,
trois $T$-groupoïdes connexes $\mathcal{G},\mathcal{G}',\mathcal{G}''$

\margin{266}

\margin{\begin{tabular}{c}
266\\
\texttt{https://grothendieck.umontpellier.fr}\\
Université de Montpellier, UMR CNRS 5149, Case courrier 051, Place Eugène Bataillon\\
34095 Montpellier cedex 5 - France
\end{tabular}}

\scanpage{268}

(savoir $\underline{\mathrm{Gr}}(\underline{\mathcal{A}}_x)$,
$\underline{\unsure{\mathrm{Revuni}}}(\unsure{\mathbb{S}}_x)$,
$\underline{\unsure{\mathrm{Revuni}}}(\unsure{\mathbb{D}}_x^*)$),
des équivalences \deleted{\unsure{de}} de $T$-groupoïdes
\[
(111)\quad
\diagram{\begin{tikzcd}[ampersand replacement=\&, column sep=large]
\mathcal{G} \arrow[r,"\substack{F\\\approx}"] \& \mathcal{G}',
\& \mathcal{G} \arrow[r,"\substack{G\\\approx}"] \& \mathcal{G}''
\end{tikzcd}}
\]
et des objets $x'\in\mathrm{Ob}\,\mathcal{G}'$,
$x''\in\mathrm{Ob}\,\mathcal{G}''$ (savoir
$\unsure{\mathbb{S}}^{\sim}$ et $\unsure{\mathbb{D}}_x^{*\sim}$),
et si on considère les foncteurs
\[
(112)\quad\left\{
\diagram{\begin{tikzcd}[ampersand replacement=\&, column sep=small,row sep=large]
\mathcal{G} \arrow[r,"f=f_{x'}"] \& \underline{\mathrm{Tors}}(T),
\& \mathcal{G} \arrow[r,"g=g_{x''}"] \& \underline{\mathrm{Tors}}(T) \\
x \arrow[r,mapsto] \& \mathrm{Hom}_{\mathcal{G}'}(F(x),x'),
\& x \arrow[r,mapsto] \& \mathrm{Hom}_{\mathcal{G}''}(G(x),x'')
\end{tikzcd}}
\right.
\]
et le foncteur correspondant
\[
(113)\quad\left\{
\diagram{\begin{tikzcd}[ampersand replacement=\&, column sep=large,row sep=large]
\mathcal{G} \arrow[r] \& \underline{\mathrm{Tors}}(T) \\
x \arrow[r,mapsto] \&
\mathrm{Hom}_{\underline{\mathrm{Tors}}(\unsure{T|G})}(f(x),g(x))
\end{tikzcd}}
\right.
\;,
\]
à prouver que ce dernier foncteur est ess\textsuperscript{t}. constant.
On en a, si \deleted{$F',G'$ désignent}
\[
(114)\quad
\diagram{\begin{tikzcd}[ampersand replacement=\&, column sep=large]
F_1\deleted{\unsure{f}}:\mathcal{G}' \arrow[r,"\approx"] \& \mathcal{G},
\& \unsure{{}'G_1|G_1}:\mathcal{G}'' \arrow[r,"\approx"] \& \mathcal{G}
\end{tikzcd}}
\]
désignent des quasi-inverses de $F,G$, des isom fonctoriels en $x$
\[
(115)\quad\left\{
\begin{array}{l}
f(x)\overset{\mathrm{def}}{=}\mathrm{Hom}_{\mathcal{G}'}(F(x),x')
\simeq\mathrm{Hom}_{\mathcal{G}}(x,F_1(x')) \\
g(x)\overset{\mathrm{def}}{=}\mathrm{Hom}_{\mathcal{G}''}(G(x),x'')
\simeq\mathrm{Hom}_{\mathcal{G}}(x,F_2(x''))
\end{array}\right.
\]
d'où \unsure{résulte|résultent} des isom de $T$-torseurs
\[
(116)\quad
\unsure{\mathrm{Isom}}(f(x),g(x))
\simeq\unsure{\mathrm{Isom}}(F_1(x'),F_2(x''))
\]
d'où l'assertion.

\margin{267}

\margin{\begin{tabular}{c}
267\\
\texttt{https://grothendieck.umontpellier.fr}\\
Université de Montpellier, UMR CNRS 5149, Case courrier 051, Place Eugène Bataillon\\
34095 Montpellier cedex 5 - France
\end{tabular}}

\scanpage{269}
\mspage{782}

\uline{Définition} On dit qu'on a une \uline{corrélation}
entre un rev. univ. $\unsure{\mathbb{S}}_x^{\sim}$ de
$\unsure{\mathbb{S}}_x$, et un autre $\unsure{\mathbb{D}}_x^{*\sim}$
de $\unsure{\mathbb{D}}_x^*$, quand on s'est donné un élément de
$H(\unsure{\mathbb{S}}_x^{\sim},\unsure{\mathbb{D}}_x^{*\sim})$.

Donc, si $A\in\mathcal{A}_x$, la donnée d'une corrélation revient à la
donnée d'un élément de $H(A)$, i.e. d'un isom.
\deleted{\unsure{de}} entre les $T_x$-torseurs
$\unsure{\mathbb{S}}_x^{\sim}(A)$ et
$\unsure{\mathbb{D}}_x^{*\sim}(\unsure{\mathrm{Ep}|E_p}(A))$.
Cette identification est compatible avec les relations d'inclusion
entre arcs $A\subset A'$.
\deleted{\unsure{D}} Les corrélations entre
$\unsure{\mathbb{S}}_x^{\unsure{*|\sim}}$ et
$\unsure{\mathbb{D}}_x^{*\sim}$ forment un torseur \unsure{sous|sur} $T_x$.
\deleted{Pour $\unsure{\mathbb{S}}_x^{\sim}$ fixé}

Soit $\underline{\mathrm{Corr}}_x(\unsure{\mathfrak{X}|X})
=\underline{\mathrm{Corr}}_x(\deleted{\unsure{\mathbb{S}}}\unsure{X})$
la catégorie des \deleted{\unsure{couples}}\unsure{triples}
$(\widetilde{\unsure{\mathbb{S}}}_x,
\widetilde{\unsure{\mathbb{D}}}_x^*,h)$ de rev. univ.
corrélés de $\unsure{\mathbb{S}}_x$ et de $\unsure{\mathbb{D}}_x^*$,
on a des \deleted{\unsure{équi|équiv.}} foncteurs
\[
\diagram{\begin{tikzcd}[ampersand replacement=\&, column sep=large,row sep=0.2ex]
\& \unsure{\mathbb{S}}_x^{\unsure{*\sim|\sim}} \\
(\widetilde{\unsure{\mathbb{S}}}_x,\unsure{\mathbb{D}}_x^{*\sim},h)
\arrow[ru] \arrow[rd] \& \\
\& \unsure{\mathbb{D}}_x^{*\sim}
\end{tikzcd}}
\]
\[
(117)\quad
\diagram{\begin{tikzcd}[ampersand replacement=\&, column sep=large,row sep=large]
\& \underline{\mathrm{Corr}}_x(X) \arrow[ld,"\approx"'] \arrow[rd,"\approx"] \& \\
\underline{\unsure{\mathrm{Revuni}}}(\unsure{\mathbb{S}}_x)
\& \& \underline{\unsure{\mathrm{Revuni}}}
(\deleted{\unsure{\mathbb{S}}}\unsure{\mathbb{D}}_x^*)
\end{tikzcd}}
\;,
\]
\unsure{qui} sont des équivalences de catégories.

\margin{268}

\margin{\begin{tabular}{c}
268\\
\texttt{https://grothendieck.umontpellier.fr}\\
Université de Montpellier, UMR CNRS 5149, Case courrier 051, Place Eugène Bataillon\\
34095 Montpellier cedex 5 - France
\end{tabular}}

\scanpage{270}

Toutes ces constructions ont été faites de façon intrinsèque assez
soigneusement pour \unsure{passer} aux iso. de germes de surfaces
différentiables
$\diagram{\begin{tikzcd}[ampersand replacement=\&, column sep=large]
(X,x) \arrow[r,"\substack{u\\\sim}"] \& (X',x')
\end{tikzcd}}$.
Supposons choisi, en plus de
$(\unsure{\mathbb{S}}_x^{\sim},\unsure{\mathbb{D}}_x^{*\sim},h_x)$
\unsure{relatif|relatifs} à $(X,x)$, aussi
$(\unsure{\mathbb{S}}_{x'}^{\sim},\unsure{a}\,
\unsure{\mathbb{D}}_{x'}^{*\sim},h'_{x'})$ \unsure{relatif|relatifs} à $(X',x')$.
\unsure{Or} l'iso. des germes étant donné, il induit donc
\[
(118)\quad\left\{
\diagram{\begin{array}{lll}
u_{\unsure{\mathbb{S}}}:&
\unsure{\mathbb{S}}_x\simeq\unsure{\mathbb{S}}_{x'}
&\text{iso. d'esp. top.}\\[0.7em]
u_{\unsure{\mathbb{D}}^*}:&
\unsure{\mathbb{D}}_x^*\xrightarrow{\sim}\unsure{\mathbb{D}}_{x'}^*
&\text{iso. de germes d'esp.}
\end{array}}
\right.
\]
\[
\diagram{\begin{array}{lll}
u_T:&T_x\simeq T_{x'}
&\begin{array}{l}
\text{iso des }\mathbb{Z}\text{-modules }\inserted{\text{top.}}\\
\unsure{\text{libres}}\text{ de rg }1
\end{array}
\end{array}}
\]
Considérons alors les deux ensembles
\[
(119)\quad\left\{
\begin{array}{l}
F(u)=\text{ens des }\unsure{\text{iso.}}\quad
\unsure{\mathbb{S}}_x^{\sim}\simeq\unsure{\mathbb{S}}_{x'}^{\sim}
\text{ au dessus de }\inserted{u_{\unsure{\mathbb{S}}}}
\\[0.6em]
\hfill\deleted{\unsure{u_{\mathbb{S}}}}
\\[0.6em]
G(u)=\text{ens des iso.}\quad
\unsure{\mathbb{D}}_x^{*\sim}\simeq\unsure{\mathbb{D}}_{x'}^{*\sim}
\text{ au dessus de }\inserted{u_{\unsure{\mathbb{D}}^*}},
\end{array}\right.
\]
ce sont l'un et l'autre des $T_x$-torseurs.
Ceci étant, je dis qu'on a un iso. can. de $T_x$-torseurs
\[
(120)\quad
\diagram{\begin{tikzcd}[ampersand replacement=\&, column sep=large]
F(u) \arrow[r,"\sim"] \& G(u)
\end{tikzcd}}
\]
--- i.e. ``il revient au $=$ de relever $u_{\unsure{\mathbb{S}}}$ en
$u_{\widetilde{\unsure{\mathbb{S}}}}:
\unsure{\mathbb{S}}_x^{\sim}\simeq S_{x'}^{\sim}$,
ou de relever $u_{\unsure{\mathbb{D}}^*}$ en
\deleted{$\unsure{u_{\widetilde{\mathbb{S}}}}$}
$u_{\unsure{\mathbb{D}}^{*\sim}}:
\unsure{\mathbb{D}}_x^{*\sim}\unsure{\longrightarrow|-}
\unsure{\mathbb{D}}_{x'}^{*\sim}$''.
Pour le voir, on note que par
\deleted{fonctorialité des \unsure{constructions}}
\inserted{transport de structure} les

\margin{269}

\margin{\begin{tabular}{c}
269\\
\texttt{https://grothendieck.umontpellier.fr}\\
Université de Montpellier, UMR CNRS 5149, Case courrier 051, Place Eugène Bataillon\\
34095 Montpellier cedex 5 - France
\end{tabular}}

\scanpage{271}
\mspage{783}

images inverses par $u_{\unsure{\mathbb{S}|\mathcal{S}}}$ et par $u_{\unsure{\mathbb{D}|D}^{*}}$ de
$\unsure{\mathbb{S}|\mathcal{S}}_{x'}^{\sim}$ et de
$\unsure{\mathbb{D}|D}_{x'}^{*\sim}$,
\unsure{soient|soit} $\unsure{\mathbb{S}|\mathcal{S}}_x^{\prime\sim}$, $\unsure{\mathbb{D}|D}_x^{*\prime\sim}$,
sont munis d'une corrélation, que je désigne par $h'_x$.
Ceci posé, \deleted{\unsure{l'existence} d'un iso
$F(u)\simeq G(u)$} des iso. de $T_x$-torseurs canoniques
\[
\begin{array}{rl}
(121)&F(u)\simeq\operatorname{Isom}_{\unsure{\mathbb{S}|\mathcal{S}}_x}
  (\unsure{\mathbb{S}|\mathcal{S}}_x^{\sim},\unsure{\mathbb{S}|\mathcal{S}}_x^{\prime\sim})\\[3pt]
&G(u)\simeq\operatorname{Isom}_{\unsure{\mathbb{D}|D}_x^{*}}
  (\unsure{\mathbb{D}|D}_x^{*\sim},\unsure{\mathbb{D}|D}_x^{*\prime\sim})
\end{array}
\]
Or par le diagramme d'équivalences (117) on a bien des iso. de
$T_x$-torseurs
\[
\begin{gathered}(122)\\[3pt]
\diagram{\begin{tikzcd}[ampersand replacement=\&, column sep=small,row sep=large,column sep=9pt,row sep=14pt,execute at begin picture={\fontsize{10}{12}\selectfont}]
\& {\begin{gathered}\operatorname{Isom}((\unsure{\mathbb{S}|\mathcal{S}}_x^{\sim},\unsure{\mathbb{D}|D}_x^{*\sim},h_x),\\[-1pt]
(\unsure{\mathbb{S}|\mathcal{S}}_x^{\prime\sim},\unsure{\mathbb{D}|D}_x^{*\prime\sim},h'_x))\end{gathered}}
\arrow[dl,"\sim"']\arrow[dr,"\sim"] \& \\
{\begin{gathered}F(u)=\\[-1pt]\operatorname{Isom}(\unsure{\mathbb{S}|\mathcal{S}}_x^{\sim},\unsure{\mathbb{S}|\mathcal{S}}_x^{\prime\sim}\unsure{]|)}\end{gathered}}
\&\&{\begin{gathered}G(u)=\\[-1pt]\unsure{\operatorname{Isom}}(\unsure{\mathbb{D}|D}_x^{*\sim},\unsure{\mathbb{D}|D}_x^{*\prime\sim})\end{gathered}}
\end{tikzcd}}
\quad ,
\end{gathered}\]
d'où la bijection canonique \unsure{cherchée} (120) ---
\deleted{\unsure{des}} qui donne en particulier (96).

\deleted{\unsure{La}} La bijection canonique (120) a la propriété
$\pm$ évidente suivante, qui la caractérise :
\[
(123)\quad\left\{\vcenter{\hbox{\begin{minipage}{0.82\linewidth}
\deleted{\unsure{On}} Soient $u_{\unsure{\mathbb{S}|\mathcal{S}}^{\sim}}$ et
$u_{\unsure{\mathbb{D}|D}^{*\sim}}$ des
\deleted{\unsure{prolongements}}\inserted{relèvements} de $u$
\deleted{\unsure{$u_{\mathbb{S}}$}}\deleted{\unsure{associés}}
$\unsure{\mathbb{S}|\mathcal{S}}_x^{\sim}\simeq\unsure{\mathbb{S}|\mathcal{S}}_{x'}^{\sim}$,
$\unsure{\mathbb{D}|D}_x^{*\sim}\longrightarrow\unsure{\mathbb{D}|D}_{x'}^{*\sim}$,
associés par (120), soit $A\in\mathcal{A}_x$,
$A'\in\mathcal{A}'_x$ son image par
$u_{\unsure{\mathbb{S}|\mathcal{S}}}$, soit $\alpha$ une section de $\unsure{\mathbb{S}|\mathcal{S}}_x^{\sim}$ sur $A$,
$\beta$ la section de $\unsure{\mathbb{D}|D}_x^{*\sim}$ sur
$\unsure{\mathrm{Ep}|E_p}(A)$ \unsure{associée|associé} par la corrélation
$h_x$, soit $\alpha'$ la section de $\unsure{\mathbb{S}|\mathcal{S}}_{x'}^{\sim}$ sur $A'$
déduite de $\alpha$ en appliquant $u_{\unsure{\mathbb{S}|\mathcal{S}}^{\sim}}$,
$\beta'$ la section de $\unsure{\mathbb{D}|D}_{x'}^{*\sim}$ sur
$\mathrm{Ep}(A')$ déduite de $\beta$ en appliquant
$u_{\unsure{\mathbb{D}|D}^{*\sim}}$, alors $\alpha'$ et $\beta'$ sont
\unsure{associés} par la corrélation $h'$.
\end{minipage}}}\right.
\]
\margin{270}

\margin{\begin{tabular}{c}
270\\
\texttt{https://grothendieck.umontpellier.fr}\\
Université de Montpellier, UMR CNRS 5149, Case courrier 051, Place Eugène Bataillon\\
34095 Montpellier cedex 5 - France
\end{tabular}}

\scanpage{272}

Ceci \deleted{\unsure{caract.}} \unsure{(}pour $A,\alpha$ fixés)
caractérise l'application \deleted{bijection} (120), car pour
$u_{\unsure{\mathbb{S}|\mathcal{S}}^{\sim}}$ fixé, il détermine $u_{\unsure{\mathbb{D}|D}^{*\sim}}$ ainsi :
on a $\beta$ par $h_x$, $\alpha'$ \unsure{déduit|déduite} de $\alpha$ par
$u_{\unsure{\mathbb{S}|\mathcal{S}}}$, $\beta'$ \unsure{déduit|déduite} de $\alpha'$ par $h_{x'}$,
or $\beta'$ doit être égal à la section déduite de $\beta$ en appliquant
$u_{\unsure{\mathbb{D}|D}^{*\sim}}:\unsure{\mathbb{D}|D}_x^{*\sim}
\unsure{\longrightarrow|-}\unsure{\mathbb{D}|D}_{x'}^{*\sim}$,
et ceci détermine bien $u_{\unsure{\mathbb{D}|D}^{*\sim}}$ de façon unique ...

\uline{NB} On peut utiliser cette construction, pour $A,\alpha$ fixés,
pour définir l'application (120), on trouve que c'est un
$T_x$-morphisme (donc iso) et,
\deleted{\unsure{qu'elle|puisqu'elle}}\inserted{du $\unsure{\overline{\mathrm m}|=}$}
coup, qu'elle ne dépend pas du choix de $\alpha$.
Mais a priori il dépend \unsure{peut-être} du choix de $A$, et le point qui
demande vérification est que cette bijection en fait ne dépend pas du choix
de $A$. On sait en effet que pour $A\subset A'$, elle est la
$\unsure{\overline{\mathrm m}|=}$ pour $A$ et $A'\unsure{;|,}$
d'où aussitôt le résultat, car la relation d'équiv. engendrée dans
$\mathcal{A}_x$ par la rel. d'inclusion est la relation grossière .... .

\margin{271}

\margin{\begin{tabular}{c}
271\\
\texttt{https://grothendieck.umontpellier.fr}\\
Université de Montpellier, UMR CNRS 5149, Case courrier 051, Place Eugène Bataillon\\
34095 Montpellier cedex 5 - France
\end{tabular}}

\scanpage{273}
\mspage{784}

F) \uline{Application} : \uline{interprétation différentielle des groupes
de Teichmüller spéciaux}.

Soit $X$ un espace topologique, muni d'une partie $I$ discrète
\inserted{finie}, telle que $X$ soit une surface au voisinage de $I$ ---
de plus on se donne une structure différentiable \inserted{$(C_1)$}
au voisinage \unsure{des|de} $I$. Soit $G$ \deleted{\illegible{1}} un groupe
topologique, opérant continûment sur $(X,I)$, et respectant la structure
différentiable de $X$ au voisinage de $S$. \unsure{Soient|Soit}
\[
(124)\quad\left\{\begin{array}{l}
\unsure{\mathbb{S}|\mathcal{S}}_{X,I}\ \unsure{\text{ou}|\sim}\ \unsure{\mathbb{S}|\mathcal{S}}_I
=\displaystyle\coprod_{x\in I}
\underbrace{\unsure{\mathbb{S}|\mathcal{S}}_{X,x}}_{\text{ou }\unsure{\mathbb{S}|\mathcal{S}}_x}\ ,\\[8pt]
\unsure{\mathbb{D}|D}_{X,I}^{*}\ \unsure{\text{ou}|\sim}\ \unsure{\mathbb{D}|D}_{I}^{*}
=\displaystyle\coprod_{x\in I}
\underbrace{\unsure{\mathbb{D}|D}_{X,x}^{*}}_{\text{ou }\unsure{\mathbb{D}|D}_x^{*}}
\end{array}\right.
\]
on suppose que l'opération de $G$ sur $\unsure{\mathbb{S}|\mathcal{S}}_S$ induite par
l'action de $G$ soit également continue.

Soit, pour $\forall\ x\in I$, $\unsure{\mathbb{D}|D}_x^{*\sim}$ un rev.
universel de $\unsure{\mathbb{D}|D}_x^{*}$, et posons
\[
(125)\quad
\unsure{\mathbb{D}|D}_{X,I}^{*\sim}\ \unsure{\text{ou}|-}\ \unsure{\mathbb{D}|D}_I^{*\sim}
=\coprod_{x\in S}\unsure{\mathbb{D}|D}_x^{*\sim}\ .
\]
Oublions pour l'instant les \unsure{données} différentiables,
et considérons le groupe
\[
(126)\quad
\widetilde{G}=\left\{u,\widetilde{u}\ \middle|\
\begin{array}{l}
u\in G,\ \widetilde{u}\text{ un relèvement de}\\
u_{\unsure{\mathbb{D}|D}_I^{*}}\text{ en un \unsure{autom.} }
 u_{\unsure{\mathbb{D}|D}_I^{*\sim}}\\
\text{de }\unsure{\mathbb{D}|D}_I^{*\sim}
\end{array}\right\}
\]
On a un hom. $(u,\widetilde{u})\longmapsto\widetilde{u}$,
$\widetilde{G}\longrightarrow G$, qui est

\margin{272}

\margin{\begin{tabular}{c}
272\\
\texttt{https://grothendieck.umontpellier.fr}\\
Université de Montpellier, UMR CNRS 5149, Case courrier 051, Place Eugène Bataillon\\
34095 Montpellier cedex 5 - France
\end{tabular}}

\scanpage{274}

épi, et, posant
\[
(127)\quad T_I=\prod_{x\in I}T_x\ ,
\]
on a une suite exacte
\[
(128)\quad
\diagram{\begin{tikzcd}[ampersand replacement=\&, column sep=small]
1\arrow[r] \& T_I\arrow[r,thick] \& \widetilde{G}\arrow[r] \& G\arrow[r] \& 1
\end{tikzcd}}\ .
\]
Je voudrais introduire sur $\widetilde{G}$ une topologie telle que (128)
devienne une suite exacte de groupes topologiques.
Quitte à remplacer $G$ par le sous-groupe $G'$ ouvert d'indice fini
noyau de l'hom canonique
\[
(129)\quad
\diagram{\begin{tikzcd}[ampersand replacement=\&]
G\arrow[r] \& \unsure{\mathfrak{S}|\mathcal{S}}_I
\end{tikzcd}}\ ,
\]
donc $\widetilde{G}$ par l'image inverse $\widetilde{G}'$ de $G'$,
on est ramené au cas où $G$ \uline{opère trivialement sur} $I$,
\deleted{et Quitte à \illegible{1}}
\inserted{Alors $\widetilde{G}$ s'exprime}
\deleted{encore à un \unsure{sous-groupe} ouvert plus petit
(d'indice divisant $2^{\operatorname{card}I}$)
\illegible{1}}
\deleted{\unsure{= que}}
comme le produit fibré sur $G$ de la famille, indexée par $I$,
des extensions
\[
(130)\quad
\diagram{\begin{tikzcd}[ampersand replacement=\&, column sep=small]
1\arrow[r] \& T_x\arrow[r] \& \widetilde{G}_x\arrow[r] \& G\arrow[r] \& 1
\end{tikzcd}}
\]
correspondant aux cas $I_x\overset{\mathrm{def}}{=}\{x\}$.
Donc on est ramené au cas où $\operatorname{card}\unsure{I|S}=1$.

\margin{273}

\margin{\begin{tabular}{c}
273\\
\texttt{https://grothendieck.umontpellier.fr}\\
Université de Montpellier, UMR CNRS 5149, Case courrier 051, Place Eugène Bataillon\\
34095 Montpellier cedex 5 - France
\end{tabular}}

\scanpage{275}
\mspage{785}

Soit \deleted{$\unsure{\Delta_x|\mathbb D_x}$}
\[
(131)\qquad
\unsure{\{}\Delta_{x,0}\deleted{\unsure{\}_{\alpha\in\mathbf N}}}
\]
\unsure{un|une} \deleted{\unsure{suite décr.}} de \unsure{voisinages|voisinage} de $x$ dans $X$ qui \unsure{soit|soient} \deleted{\unsure{des}} \inserted{un} \unsure{disques ouverts|disque ouvert}, \deleted{et que \unsure{forment} une \unsure{base} du filtre des voisinages de $x$ dans $X$, et soit}
On peut \unsure{alors} trouver un rev. universel $\Delta_{x,0}^{*\sim}$ de $\Delta_{x,0}^{*}=\Delta_{x,0}\setminus\{x\}$ qui \unsure{induise|induit} $\unsure{\mathbb D|D}_x^{*\sim}$ sur $\unsure{\mathbb D|D}_x^{*}$. Par hypothèse de continuité, il existe un voisinage $U$ de $1$ dans $G$, et un voisinage $\Delta_{x,1}$ de $x$, tel que l'on ait
\[
(132)\qquad
\diagram{\begin{tikzcd}[ampersand replacement=\&]
U\times\Delta_{x,1} \arrow[r] \arrow[d, phantom, "\cap" description] \& \Delta_{x,0} \arrow[d, phantom, "\cap" description] \\
G\times X \arrow[r] \& X
\end{tikzcd}},
\]
et on \unsure{peut} supposer que $\Delta_{x,1}$ soit un \unsure{disque ouvert|disques ouverts}. Ceci pris; pour $u\in\unsure{U}$ fixé, la donnée du relèvement $\widetilde u$ de $u_{\unsure{\mathbb D|D}_x}$ à $\unsure{\mathbb D|D}_x^{\sim}$ équivaut à celle d'un relèvement en \deleted{$\unsure{\Delta|\mathbb D}$}
\[
(133)\qquad
\diagram{\begin{tikzcd}[ampersand replacement=\&]
\widetilde\Delta_{x,1} \arrow[r, "\widetilde u_{0,1}"] \& \Delta_{x,0}^{\sim}
\end{tikzcd}}
\]
compatible avec
\[
u\mid\Delta_{x,1}^{*}:\quad
\diagram{\begin{tikzcd}[ampersand replacement=\&]
\Delta_{x,\unsure{1|0}}^{*}\arrow[r] \& \Delta_{x,0}^{*}
\end{tikzcd}}.
\]
Ceci dit, dans l'image inverse $\widetilde U$ dans $\widetilde G$ du voisinage $U$ de $1$ dans $G$, on va

\margin{274}

\margin{\begin{tabular}{c}
274\\
\texttt{https://grothendieck.umontpellier.fr}\\
Université de Montpellier, UMR CNRS 5149, Case courrier 051, Place Eugène Bataillon\\
34095 Montpellier cedex 5 - France
\end{tabular}}

\scanpage{276}

introduire la top. \inserted{la moins fine rendant continues} \deleted{de la \unsure{conv. compacte} \unsure{des|de} $\widetilde u$ et de $u$} les applications
\[
(134)\quad
\left\{\begin{array}{ll}
\diagram{\begin{tikzcd}[ampersand replacement=\&]
\unsure{p_1|p_i}:\widetilde U \arrow[r] \& U
\end{tikzcd}}
& \text{\unsure{induit|induite} par }\diagram{\begin{tikzcd}[ampersand replacement=\&]\widetilde G\arrow[r]\&G\end{tikzcd}}
\\[1.2ex]
\diagram{\begin{tikzcd}[ampersand replacement=\&]
q:\widetilde U \arrow[r, phantom, "\unsure{\Longrightarrow|\longrightarrow}" description] \arrow[r, phantom, "\deleted{\unsure{\longrightarrow}}"', shift right=2ex] \&[35pt] \operatorname{Cont}(\Delta_{x,1}^{*\sim},\Delta_{x,0}^{*\sim})
\end{tikzcd}}
& \diagram{\begin{tikzcd}[ampersand replacement=\&]
(u,\widetilde u) \arrow[r, mapsto] \& \deleted{\unsure{(\widetilde u)_{0,1}}}\inserted{\widetilde u_{0,1}}
\end{tikzcd}}
\\[-.2ex]
& \diagram{\begin{tikzcd}[ampersand replacement=\&,nodes in empty cells]
\deleted{\unsure{(u,\widetilde u)}} \arrow[r, phantom, "\deleted{\mapsto}" description] \& \deleted{(\widetilde u^{-1})_{0,1}}
\end{tikzcd}}
\\[1.2ex]
q'=q\circ\sigma,\quad\text{où }\sigma(u,\widetilde u)=(u,\widetilde u)^{-1}: & \diagram{\begin{tikzcd}[ampersand replacement=\&]
\widetilde U\arrow[r]\&\operatorname{Cont}(\Delta_{x,1}^{*\sim},\Delta_{x,0}^{*\sim})
\end{tikzcd}}
\end{array}\right.
\]

\unsure{On aura|J'aurai} pris soin de prendre $\unsure{U|\widetilde U}$ stable par $u\mapsto u^{-1}$, ce qui implique qu'il en est de $\unsure{\overline{\mathrm m}|=}$ de $\widetilde U$.

On trouve donc, à l'aide de cette topologie sur $\widetilde U$, \deleted{\unsure{un} \illegible{1} \unsure{un filtre}} le filtre des voisinages de $1_{\widetilde G}$ dans $\widetilde U$, donc aussi le filtre \unsure{image} de celui-ci dans $\widetilde G$. Il faut vérifier que ce dernier est bien le filtre des voisinages de $1$ pour une topologie de groupe top., et que pour celle-ci, $\widetilde G\to G$ est un rev. de $G$. \deleted{\unsure{Je suppose}}\inserted{J'admets} que tout marche.

Quand $X$ est \inserted{une surface} compacte orientable, \inserted{connexe}, et \deleted{que} si $\unsure{T|\overline T}$ désigne son module des

\margin{275}

\margin{\begin{tabular}{c}
275\\
\texttt{https://grothendieck.umontpellier.fr}\\
Université de Montpellier, UMR CNRS 5149, Case courrier 051, Place Eugène Bataillon\\
34095 Montpellier cedex 5 - France
\end{tabular}}

\scanpage{277}
\mspage{786}

orientations, alors tous les $T_x$ sont can. isom. à $\unsure{T|\overline T}$
\[
(135)\qquad T_x\simeq\unsure{T|\overline T}\qquad\forall x\in I,
\]
et la suite exacte (128) prend la forme
\[
(136)\qquad
\diagram{\begin{tikzcd}[ampersand replacement=\&]
1\arrow[r, dash]\&\unsure{T^I|\overline T^I}\arrow[r]\&\widetilde G\arrow[r]\&G\arrow[r]\&1
\end{tikzcd}}
\]
Le cas ``universel'' est celui où
\[
(137)\qquad G=\mathcal A=\operatorname{Aut\,top}(X,S)
\]
(avec la top. de la conv. uniforme des $u$ et des $u^{-1}$), Alors dans les cas \uline{\unsure{anabéliens}} (type de $(X,S)$ étant $g,\nu$, et
\[
(138)\qquad\begin{array}{l}
2g-2+\nu>0\quad\text{i.e.}\quad 2g+\nu\geq3\\[.5ex]
\text{i.e.}\quad(g,\nu)\ne(0,0),(0,1),(0,2),(1,0)
\end{array}
\]
\unsure{)|,} alors on sait que $\mathcal A^0$ est \unsure{1-connexe} (\unsure{est|et} $\unsure{\overline{\mathrm m}|=}$ contractile) donc \unsure{tte|toute} \deleted{\unsure{revêtement|revêtement de}}\inserted{extension} de $\mathcal A^0$ par un groupe discret est canoniquement scindée, et l'extension $\widetilde{\mathcal A}$ de $G$ provient alors d'une extension $\widetilde{\unsure{\mathcal T|\tau}}$
\[
(139)\qquad
\diagram{\begin{tikzcd}[ampersand replacement=\&]
1\arrow[r]\&\unsure{\mathbb Z^I|T^I}\arrow[r, dash]\&\widetilde{\unsure{\mathcal T|\tau}}\arrow[r]\&\unsure{\mathcal T|\tau}\arrow[r]\&1
\end{tikzcd}}
\]
du groupe de Teichmüller
\[
(140)\qquad
\unsure{\mathcal T|\tau}=\unsure{\mathcal T|\tau}(X,S)=\mathcal A/\mathcal A^0=\pi_0\bigl(\underbrace{\uline{\operatorname{Aut\,top}}(X,S)}_{\unsure{A|\mathcal A}}\bigr)
\]

\margin{276}

\margin{\begin{tabular}{c}
276\\
\texttt{https://grothendieck.umontpellier.fr}\\
Université de Montpellier, UMR CNRS 5149, Case courrier 051, Place Eugène Bataillon\\
34095 Montpellier cedex 5 - France
\end{tabular}}

\scanpage{278}

\unsure{qui|que} prend ainsi le nom de \uline{groupe de
Teichmüller spécial} de $(X,S)$. Pour
$G$ quelconque l'action de $G$ sur
$(X,S)$ \unsure{revient} à la donnée d'un hom.
de groupes topologiques
\diagram{\begin{tikzcd}[ampersand replacement=\&]
G \arrow[r] \& \unsure{\mathcal A|A}
\end{tikzcd}}
\unsure{s'insérant} dans un diagramme de
groupes top.
\[
\tag{141}
\diagram{\begin{tikzcd}[ampersand replacement=\&]
G \arrow[r] \arrow[d] \& \unsure{\mathcal A|A} \arrow[d] \\
\pi_0(G) \arrow[r] \& \unsure{\mathcal T|\tau}=\pi_0(\unsure{\mathcal A|A})
\end{tikzcd}},
\]
et on trouve alors, par fonctorialité,
que le $\widetilde G$ construit précédemment
est l'image inverse de \inserted{l'ext.}
$\widetilde{\unsure{\mathcal T|\tau}}=\operatorname{Sp}\unsure{\mathcal T|\tau}$
$\unsure{a|\alpha}$ de $\unsure{\mathcal T|\tau}$
par $\mathbb Z^I$, via $G\longrightarrow\unsure{\mathcal T|\tau}$.

\unsure{Ces} préliminaires étant \unsure{pris},
revenons à la situation où $X$
admet une structure diff. $C_1$ au voisinage
de $\unsure{I|S}$, et soit pour $\forall\,x\in\unsure{I|S}$,

\margin{277}

\margin{\begin{tabular}{c}
277\\
\texttt{https://grothendieck.umontpellier.fr}\\
Université de Montpellier, UMR CNRS 5149, Case courrier 051, Place Eugène Bataillon\\
34095 Montpellier cedex 5 - France
\end{tabular}}

\scanpage{279}
\mspage{787}

$\widetilde{\unsure{\mathbb S|\mathcal S}}_x$
``\unsure{le}'' rev. universel de $\unsure{\mathbb S|\mathcal S}_x$ \unsure{corrélé} au
rev. universel $\unsure{\mathbb D|D}_x^{*\sim}$ de $\unsure{\mathbb D|D}_x^*$. Posons
\[
\tag{142}
\widetilde{\unsure{\mathbb S|\mathcal S}}_I
=\coprod_{x\in I}\widetilde{\unsure{\mathbb S|\mathcal S}}_x,
\]
et notons que l'on a, pour $\forall\,u\in G$,
une correspondance biunivoque entre
l'ens. des relèvements $\widetilde u_{\unsure{\mathbb D|D}_I^*}$ de $u_{\unsure{\mathbb D|D}_I^*}$
en autom. de $\unsure{\mathbb D|D}_I^{*\sim}$, et l'ens. des
relèvements $\widetilde u_{\unsure{\mathbb S|\mathcal S}_I}$ de
$u_{\unsure{\mathbb S|\mathcal S}}$ en autom. de
$\widetilde{\unsure{\mathbb S|\mathcal S}}_I$,
qui soit un isom. de $T_I$-torseurs.
Cette correspondance est compatible avec
les compositions, et on trouve ainsi
un isom. canonique
\[
\tag{143}
\widetilde G\simeq\text{groupe des }
\left\{u,\widetilde u_{\unsure{\mathbb S|\mathcal S}}\;\middle|\;
\begin{array}{l}
u\in G,\ \widetilde u_{\unsure{\mathbb S|\mathcal S}_I}\ \text{rel.}\\
\text{de }u_{\unsure{\mathbb S|\mathcal S}_I}\ \text{en un}\\
\text{autom. de }\unsure{\mathbb S|\mathcal S}_I^{\sim}
\end{array}\right\}
\]

On vérifie que c'est là un \uline{isom.}
\uline{top.} en mettant sur le deuxième
\deleted{\unsure{met}} la top. la moins fine rendant
continues les applications
\[
\tag{144}
\begin{array}{l}
\diagram{\begin{tikzcd}[ampersand replacement=\&, column sep=large]
\unsure{p_1^{a}|p_1}:\widetilde G \arrow[r] \& G
\end{tikzcd}}\\
\diagram{\begin{tikzcd}[ampersand replacement=\&, column sep=large]
q_1:\widetilde G \arrow[r] \&
\operatorname{Aut\,top}(\widetilde{\unsure{\mathbb S|\mathcal S}}_I)
\end{tikzcd}}
\qquad (u,\widetilde u_{\unsure{\mathbb S|\mathcal S}})
\longmapsto\widetilde u_{\unsure{\mathbb S|\mathcal S}}\\
\diagram{\begin{tikzcd}[ampersand replacement=\&, column sep=large]
q'_1=q_1\circ\sigma:\widetilde G \arrow[r] \&
\operatorname{Aut\,top}(\widetilde{\unsure{\mathbb S|\mathcal S}}_I)
\end{tikzcd}}
\end{array}
\]
\margin{\unsure{trans|tras}}

\margin{278}

\margin{\begin{tabular}{c}
278\\
\texttt{https://grothendieck.umontpellier.fr}\\
Université de Montpellier, UMR CNRS 5149, Case courrier 051, Place Eugène Bataillon\\
34095 Montpellier cedex 5 - France
\end{tabular}}

\scanpage{280}

où $\operatorname{Aut\,top}(\widetilde{\unsure{\mathbb S|\mathcal S}}_I)$ est muni de la
convergence compacte.

Dans le cas où $X$ est \inserted{une surface} comp. et diff.
orientable, on trouve donc une interprétation
des groupes précédents,
\deleted{\unsure{construit|explicite}, \unsure{par}}
\unsure{voire|voie} différentielle,
\inserted{en \unsure{tt} \unsure{qu'extension|q'extension} de $\unsure{G|\widetilde G}$ par $\unsure{T_I|\overline T_I}$,}
comme image
inverse d'un groupe de Teichmüller
spécial $\unsure{S\mathcal T|\operatorname{Sp}\mathcal T}\longrightarrow\unsure{\mathcal T|\tau}$,
par l'hom. canonique
$G\longrightarrow\unsure{\mathcal T|\tau}$ de (141).

\margin{279}

\margin{\begin{tabular}{c}
279\\
\texttt{https://grothendieck.umontpellier.fr}\\
Université de Montpellier, UMR CNRS 5149, Case courrier 051, Place Eugène Bataillon\\
34095 Montpellier cedex 5 - France
\end{tabular}}
